\section{An infinite series of arrangements of \texorpdfstring{$6\cdot2^t+2$}{6·2ᵗ + 2} straight lines}\label{sec:series}

We now construct an infinite series of arrangements of an even number of lines.
The starting arrangement of the series has $N_0=8$ lines and two triple points.
The other arrangements are obtained from it by temporarily removing one line, applying the Bartholdi--Blanc--Loisel doubling~\cite{bbl} the required number of times, and restoring the removed line.

In the notation of~\cite{bbl}, $n$ has a special meaning: it is the number of added lines.
To avoid confusion, in this section, unlike the others, the number of lines in an arrangement is always denoted by $N$.
An arrangement $\mathcal A$ of $N$ straight lines in the affine plane is a set of $N$ distinct lines, no two of which are parallel.
The \emph{multiplicity} of a point is the number of lines of the arrangement passing through it.
A point of multiplicity 2 is called \emph{simple}, and a point of higher multiplicity is called \emph{multiple}.
An arrangement is \emph{simple} if all its intersection points are simple.

The lines divide the plane into \emph{faces}.
A face is \emph{triangular} if it is bounded and its sides lie on exactly three lines of the arrangement.
We denote by $a_3(\mathcal A)$ the number of triangular faces of the arrangement $\mathcal A$.

\begin{theorem}\label{thm:series}
For every integer $t\ge0$ there exists an arrangement $\mathcal{C}_t$ of
\begin{equation}\label{eq:N-t}
N_t=6\cdot2^t+2
\end{equation}
lines with two triple points and exactly
\begin{equation}\label{eq:a3-series}
a_3(\mathcal{C}_t)=\frac{N_t(N_t-\frac52)+1}{3}
\end{equation}
triangular faces.
\end{theorem}

\subsection{The starting arrangement}

Fix an arbitrary integer $t\ge0$ for which the theorem is to be proved. For the intermediate values $s$, $0\le s\le t$, set
\[
n_s=6\cdot2^s.
\]
Consider the eight lines given by the equations
\begin{equation}\label{eq:base-lines}
\begin{array}{ll}
R:y=x-\delta,
&Y_0:y=0,\\[4pt]
L_1:y=-\dfrac12\left(x+\sqrt3\right),
&L_2:y=-\eta\Bigl(x+\dfrac1{\sqrt3}\Bigr),\\[6pt]
L_3:y=\eta\Bigl(x-\dfrac1{\sqrt3}\Bigr),
&L_4:y=\dfrac25\left(x-\sqrt3\right),\\[8pt]
L_5:y=u(\varepsilon,\eta)(x+\varepsilon),
&L_6:y=v(\varepsilon)(x-\varepsilon).
\end{array}
\end{equation}
These equations have three free parameters. We tie them to the target value $t$:
\begin{equation}\label{eq:params}
\varepsilon=\frac1{4n_t},
\qquad
\eta=3\sqrt3\,n_t^6,
\qquad
\delta=2n_t^6.
\end{equation}
The remaining two dependent parameters $u(\varepsilon,\eta)$ and $v(\varepsilon)$ are defined as follows:
\begin{equation}\label{eq:star-slopes}
u(\varepsilon,\eta)=-\frac{4\eta}{2\eta-3+\varepsilon\sqrt3(2\eta+1)},
\qquad
v(\varepsilon)=\frac{4\sqrt3}{\sqrt3+9\varepsilon}.
\end{equation}

With these equations, the eight lines form an arrangement with 15 triangular faces (\cref{fig:seed}).
The line $Y_0$ is horizontal, and the lines $L_1,\ldots,L_6$ intersect it at the points $x=\pm\varepsilon,\tan(\pm\pi/6),\tan(\pm2\pi/6)$,
as the doubling requires~\cite[Prop.~3.1]{bbl}. The parameter $\varepsilon$ satisfies the condition under which the doubling can be iterated~\cite[Remark~3.2]{bbl}.
The slopes of $L_1,L_4$ are fixed and equal to $-1/2, 2/5$, respectively. The slopes of $L_2,L_3$ are set by \eqref{eq:params} and depend on $t$.
The slopes \eqref{eq:star-slopes} of $L_5,L_6$ are chosen so that three lines meet at each of two points:
\[
P=L_1\cap L_4\cap L_6,
\qquad
Q=L_1\cap L_3\cap L_5.
\]
The parameters $\eta,\delta$ ensure that after $t$ iterations of the doubling all new intersection points lie on one side of the line $R$.
This line plays a special role:
removing it before the doubling makes the doubling possible, and afterwards the line is restored.
We call a triangular face with a side on $R$ an \emph{$R$-triangle}, and write $r(\cdot)$ for the number of $R$-triangles of an arrangement.

All lines in \eqref{eq:base-lines} have finite slope, that is, none of them is vertical.
Therefore the order of the intersection points along any line of the arrangement is uniquely determined by the order of their $x$-coordinates, and below the words ``to the right'' and ``rightmost'' refer to comparing $x$-coordinates.

\begin{figure}[t]
\centering
\resizebox{\ifdim\width>\linewidth\linewidth\else\width\fi}{!}{\begin{tikzpicture}[x=1.9486cm,y=1.9486cm]
  \begin{scope}
    \clip (-2.95083,-2.95083) rectangle (2.95083,2.95083);
    \fill[rface] (1.77671,0.4123) -- (1.95534,0.59094) -- (1.50876,0.4123) -- cycle;
    \fill[rface] (0.53269,-0.83171) -- (0.33421,-1.03019) -- (0.15497,-0.64285) -- cycle;
    \fill[rface] (-0.39154,-1.75594) -- (0.02321,-1.34119) -- (-0.22329,-2.64765) -- cycle;
    \fill[rface] (-0.39154,-1.75594) -- (-1.33642,-2.70083) -- (-0.57764,-0.76958) -- cycle;
    \fill[rface] (0.02321,-1.34119) -- (0.33421,-1.03019) -- (0.15497,-0.64285) -- cycle;
    \fill[face] (-1.95534,0.4123) -- (-0.80064,0.4123) -- (-0.68036,-0.22519) -- cycle;
    \fill[face] (-0.80064,0.4123) -- (-0.33329,0.4123) -- (-1.12239,2.11758) -- cycle;
    \fill[face] (0.35406,0.4123) -- (1.50876,0.4123) -- (0.2598,-0.08728) -- cycle;
    \fill[face] (0.35406,0.4123) -- (-0.11329,0.4123) -- (0.78586,2.70083) -- cycle;
    \fill[face] (-0.33329,0.4123) -- (-0.11329,0.4123) -- (-0.21431,0.15519) -- cycle;
    \fill[face] (-0.68036,-0.22519) -- (-0.41574,-0.3575) -- (-0.63858,-0.44663) -- cycle;
    \fill[face] (-0.41574,-0.3575) -- (0.15497,-0.64285) -- (-0.04559,-0.20943) -- cycle;
    \fill[face] (-0.63858,-0.44663) -- (-0.57764,-0.76958) -- (-0.41574,-0.3575) -- cycle;
    \fill[face] (0.2598,-0.08728) -- (0.15497,-0.64285) -- (-0.04559,-0.20943) -- cycle;
    \fill[face] (-0.04559,-0.20943) -- (-0.41574,-0.3575) -- (-0.21431,0.15519) -- cycle;
    \draw[rline] (-1.58642,-2.95083) -- (2.95083,1.58642);
    \draw[oldline] (-2.95083,0.4123) -- (2.95083,0.4123);
    \draw[oldline] (-2.95083,0.91005) -- (2.95083,-2.04078);
    \draw[oldline] (-1.27961,2.95083) -- (-0.16609,-2.95083);
    \draw[oldline] (-0.28049,-2.95083) -- (0.83303,2.95083);
    \draw[oldline] (-2.95083,-1.37153) -- (2.95083,0.98913);
    \draw[oldline] (-1.50797,2.95083) -- (1.22297,-2.95083);
    \draw[oldline] (-1.43465,-2.95083) -- (0.88408,2.95083);
    \fill[starpt] (-0.41574,-0.3575) circle[radius=1.6pt];
    \fill[starpt] (0.15497,-0.64285) circle[radius=1.6pt];
  \end{scope}
  \node[starlbl,anchor=291] at (-0.41574,-0.3575) {$P$};
  \node[starlbl,anchor=206] at (0.15497,-0.64285) {$Q$};
  \node[rimlbl,anchor=west] at (2.95083,-2.04078) {$L_{1}$};
  \node[rimlbl,anchor=west] at (2.95083,0.4123) {$Y_0$};
  \node[rimlbl,anchor=west] at (2.95083,0.98913) {$L_{4}$};
  \node[rimlbl,anchor=west] at (2.95083,1.58642) {$R$};
  \node[rimlbl,anchor=south] at (-1.54389,2.95083) {$L_{5}$};
  \node[rimlbl,anchor=south] at (-1.24368,2.95083) {$L_{2}$};
  \node[rimlbl,anchor=south] at (0.70845,2.95083) {$L_{3}$};
  \node[rimlbl,anchor=south] at (1.00866,2.95083) {$L_{6}$};
  \node[rimlbl,anchor=north] at (-1.66064,-2.95083) {$R$};
  \node[rimlbl,anchor=north] at (-1.36043,-2.95083) {$L_{6}$};
  \node[rimlbl,anchor=north] at (-0.3734,-2.95083) {$L_{3}$};
  \node[rimlbl,anchor=north] at (-0.07318,-2.95083) {$L_{2}$};
  \node[rimlbl,anchor=north] at (1.22297,-2.95083) {$L_{5}$};
  \node[rimlbl,anchor=east] at (-2.95083,-1.37153) {$L_{4}$};
  \node[rimlbl,anchor=east] at (-2.95083,0.4123) {$Y_0$};
  \node[rimlbl,anchor=east] at (-2.95083,0.91005) {$L_{1}$};
\end{tikzpicture}
}
\caption{The starting arrangement \eqref{eq:base-lines}, with 15 triangular faces (5 of them $R$-triangles, hatched) and two triple points $P$ and $Q$. The values of $\varepsilon,\eta,\delta$ differ from \eqref{eq:params} for readability, but the order of intersection points on each line is the same.}
\label{fig:seed}
\end{figure}

\begin{lemma}\label{lem:seed}
The starting arrangement has exactly 15 triangular faces, and 5 of them are $R$-triangles.
After the line $R$ is removed, the arrangement
\[
\mathcal A_0=\{Y_0,L_1,\ldots,L_6\}
\]
has exactly two triple points $P$ and $Q$ (and no other multiple points) and 10 triangular faces.
\end{lemma}

\begin{proof}
\Cref{fig:seed} illustrates the statement, which is verified in \cref{app:seed-faces}.
\end{proof}

We construct the arrangements $\mathcal A_1,\ldots,\mathcal A_t$ recursively. The arrangement $\mathcal A_0$ has been defined above.
The arrangement $\mathcal A_{s+1}$ is obtained from $\mathcal A_s$ by adding the $n_s$ lines $M_j$ given by the conditions of~\cite[Prop.~3.1]{bbl} with $n=n_s$.
By construction, $\mathcal A_s$ consists of $n_s+1$ lines.
The line $R$ is not used in the doublings; after $t$ doublings, we add it to $\mathcal A_t$. This gives the required arrangement $\mathcal C_t$:
\[
\mathcal C_t=\mathcal A_t\cup\{R\},
\]
containing $N_t=n_t+2$ lines \eqref{eq:N-t}, for which we prove the statement of the theorem.

\subsection{The doubling step}

Let $n$ be even, and let the arrangement
\[
\mathcal A=\{Y_0,L_1,\ldots,L_n\}
\]
be obtained at one of the steps of the construction. Choose coordinates in which
\[
Y_0:y=0,
\qquad
L_i:y=m_i(x-a_i).
\]
The added lines are those of~\cite[Prop.~3.1]{bbl}:
\begin{equation}\label{eq:new-lines}
\begin{gathered}
M_j:y=\mu_j\left(x-b_j\right),
\qquad
b_j=\tan\beta_j,
\qquad
\beta_j=-\frac\pi2+\left(j-\frac12\right)\frac\pi n,\\
\mu_j=-\frac{m_{\min}}{n^{10}}
\left(\sin(2\beta_j)+\frac1{n^6b_j}\right),
\qquad
m_{\min}=\min_{1\le i\le n}|m_i|,
\qquad
j=1,\ldots,n.
\end{gathered}
\end{equation}
The factor $-1$ in $\mu_j$ is the sign $\sigma$ of~\cite[Prop.~3.1]{bbl}, which equals $-1$ when $L_{n-1}$ and $L_n$ meet below $Y_0$, as $L_5$ and $L_6$ in \eqref{eq:base-lines} do.
The result of the first doubling step is shown in \cref{fig:double}.
From now on, the indices $i$ and $j$ always run over $1,\ldots,n$. The subscript of $L_i$, $a_i$, $m_i$ refers to an old line, and the subscript of $M_j$, $b_j$, $\beta_j$, $\mu_j$ refers to a new one.
Below we compare $m_{\min}$ for different arrangements, so where confusion is possible, we add the arrangement as an argument: $m_{\min}(\mathcal A)$.

\begin{figure}[t]
\centering
\resizebox{\ifdim\width>\linewidth\linewidth\else\width\fi}{!}{\begin{tikzpicture}[x=1.3379cm,y=1.3379cm]
  \begin{scope}
    \clip (-4.85853,-4.13463) rectangle (4.85853,4.13463);
    \fill[rface] (3.20474,1.93157) -- (3.19669,1.92352) -- (3.09906,1.93157) -- cycle;
    \fill[rface] (1.34815,0.07498) -- (0.5899,-0.68327) -- (-0.16484,0.83147) -- cycle;
    \fill[rface] (-0.81164,-2.08481) -- (-0.41496,-1.68813) -- (-0.63299,-3.88451) -- cycle;
    \fill[rface] (-0.81164,-2.08481) -- (-2.61147,-3.88463) -- (-1.06635,0.48108) -- cycle;
    \fill[rface] (-0.41496,-1.68813) -- (0.5899,-0.68327) -- (-0.16484,0.83147) -- cycle;
    \fill[rface] (4.60853,3.33536) -- (4.45584,3.18268) -- (3.9849,3.08591) -- cycle;
    \fill[rface] (3.943,2.66983) -- (3.95406,2.68089) -- (3.80024,2.65715) -- cycle;
    \fill[rface] (2.73951,1.46635) -- (2.71748,1.44431) -- (2.27646,1.53573) -- cycle;
    \fill[face] (-2.36504,1.93157) -- (-4.36504,1.93157) -- (-2.66685,2.08248) -- cycle;
    \fill[face] (-2.36504,1.93157) -- (-1.63299,1.93157) -- (-2.15183,1.82496) -- cycle;
    \fill[face] (-1.21034,1.93157) -- (-1.63299,1.93157) -- (-1.21879,2.01668) -- cycle;
    \fill[face] (-1.21034,1.93157) -- (-0.90094,1.93157) -- (-1.20567,1.88454) -- cycle;
    \fill[face] (-0.05564,1.93157) -- (-0.36504,1.93157) -- (-0.06017,1.88589) -- cycle;
    \fill[face] (-0.05564,1.93157) -- (0.36701,1.93157) -- (-0.04712,2.01742) -- cycle;
    \fill[face] (1.09906,1.93157) -- (0.36701,1.93157) -- (0.84917,1.83162) -- cycle;
    \fill[face] (1.09906,1.93157) -- (3.09906,1.93157) -- (1.44104,2.06836) -- cycle;
    \fill[face] (-0.71299,1.93157) -- (-0.55299,1.93157) -- (-0.61944,1.74382) -- cycle;
    \fill[face] (-0.71299,1.93157) -- (-0.90094,1.93157) -- (-0.72641,1.95851) -- cycle;
    \fill[face] (-0.55299,1.93157) -- (-0.36504,1.93157) -- (-0.54353,1.95832) -- cycle;
    \fill[face] (-1.15003,1.32407) -- (-0.82544,1.16177) -- (-1.12214,1.04309) -- cycle;
    \fill[face] (-1.15003,1.32407) -- (-2.0197,1.7589) -- (-1.20567,1.88454) -- cycle;
    \fill[face] (-0.82544,1.16177) -- (-0.16484,0.83147) -- (-0.41185,1.32721) -- cycle;
    \fill[face] (-2.66685,2.08248) -- (-3.22098,2.35954) -- (-1.85404,2.1547) -- cycle;
    \fill[face] (-2.15183,1.82496) -- (-2.0197,1.7589) -- (-3.84247,1.47758) -- cycle;
    \fill[face] (-3.22098,2.35954) -- (-3.44482,2.47146) -- (-4.60853,2.56747) -- cycle;
    \fill[face] (-4.30003,2.89907) -- (-3.44482,2.47146) -- (-1.43908,2.30598) -- cycle;
    \fill[face] (-1.12214,1.04309) -- (-1.06635,0.48108) -- (-0.82544,1.16177) -- cycle;
    \fill[face] (-1.33408,3.17807) -- (-1.24592,2.29005) -- (-0.87641,2.25956) -- cycle;
    \fill[face] (-1.23792,2.20945) -- (-1.2231,2.06016) -- (-1.85404,2.1547) -- cycle;
    \fill[face] (-1.23792,2.20945) -- (-1.24348,2.26543) -- (-1.05279,2.2259) -- cycle;
    \fill[face] (-1.21879,2.01668) -- (-1.2231,2.06016) -- (-1.09826,2.04145) -- cycle;
    \fill[face] (-1.24348,2.26543) -- (-1.24592,2.29005) -- (-1.43908,2.30598) -- cycle;
    \fill[face] (-0.10339,1.45059) -- (-0.16484,0.83147) -- (-0.41185,1.32721) -- cycle;
    \fill[face] (-0.10339,1.45059) -- (-0.06017,1.88589) -- (0.70005,1.77197) -- cycle;
    \fill[face] (0.13824,3.88463) -- (-0.01729,2.31792) -- (-0.42921,2.28131) -- cycle;
    \fill[face] (-0.01729,2.31792) -- (-0.0228,2.26243) -- (0.44886,2.35934) -- cycle;
    \fill[face] (-0.0228,2.26243) -- (-0.03001,2.18973) -- (-0.2773,2.21013) -- cycle;
    \fill[face] (-0.04249,2.06406) -- (-0.04712,2.01742) -- (-0.17411,2.04375) -- cycle;
    \fill[face] (-0.04249,2.06406) -- (-0.03001,2.18973) -- (0.49245,2.14663) -- cycle;
    \fill[face] (-0.41185,1.32721) -- (-0.82544,1.16177) -- (-0.61944,1.74382) -- cycle;
    \fill[face] (2.6596,2.55579) -- (3.9849,3.08591) -- (0.44886,2.35934) -- cycle;
    \fill[face] (2.6596,2.55579) -- (2.35558,2.43418) -- (3.80024,2.65715) -- cycle;
    \fill[face] (2.35558,2.43418) -- (1.44104,2.06836) -- (0.49245,2.14663) -- cycle;
    \fill[face] (0.70005,1.77197) -- (0.84917,1.83162) -- (2.27646,1.53573) -- cycle;
    \fill[face] (-0.86783,2.24234) -- (-0.83739,2.18125) -- (-1.05279,2.2259) -- cycle;
    \fill[face] (-0.86783,2.24234) -- (-0.87641,2.25956) -- (-0.77146,2.2509) -- cycle;
    \fill[face] (-0.79843,2.10305) -- (-0.74106,1.98792) -- (-1.09826,2.04145) -- cycle;
    \fill[face] (-0.79843,2.10305) -- (-0.83739,2.18125) -- (-0.62856,2.13796) -- cycle;
    \fill[face] (-0.72641,1.95851) -- (-0.74106,1.98792) -- (-0.63694,1.97232) -- cycle;
    \fill[face] (-0.42921,2.28131) -- (-0.44938,2.22433) -- (-0.77146,2.2509) -- cycle;
    \fill[face] (-0.46829,2.17089) -- (-0.49011,2.10926) -- (-0.62856,2.13796) -- cycle;
    \fill[face] (-0.46829,2.17089) -- (-0.44938,2.22433) -- (-0.2773,2.21013) -- cycle;
    \fill[face] (-0.53289,1.98838) -- (-0.54353,1.95832) -- (-0.63694,1.97232) -- cycle;
    \fill[face] (-0.53289,1.98838) -- (-0.49011,2.10926) -- (-0.17411,2.04375) -- cycle;
    \draw[rline] (-2.86147,-4.13463) -- (4.85853,3.58536);
    \draw[oldline] (-4.85853,1.93157) -- (4.85853,1.93157);
    \draw[oldline] (-4.85853,3.17832) -- (4.85853,-1.68021);
    \draw[oldline] (-1.42903,4.13463) -- (-0.60816,-4.13463);
    \draw[oldline] (-0.65782,-4.13463) -- (0.16305,4.13463);
    \draw[oldline] (-4.85853,-0.45146) -- (4.85853,3.43536);
    \draw[oldline] (-1.8107,4.13463) -- (2.30961,-4.13463);
    \draw[oldline] (-2.69995,-4.13463) -- (0.22672,4.13463);
    \draw[newline] (-4.85853,1.88772) -- (4.85853,2.75119);
    \draw[newline] (-4.85853,1.2688) -- (4.85853,3.26542);
    \draw[newline] (-4.85853,1.32076) -- (4.85853,2.82048);
    \draw[newline] (-4.85853,2.60493) -- (4.85853,1.1488);
    \draw[newline] (-4.85853,3.01485) -- (4.85853,1.00046);
    \draw[newline] (-4.85853,2.5881) -- (4.85853,1.78641);
    \fill[starpt] (-0.82544,1.16177) circle[radius=1.6pt];
    \fill[starpt] (-0.16484,0.83147) circle[radius=1.6pt];
  \end{scope}
  \node[lbl,anchor=west] at (4.85853,-1.68021) {$L_{1}$};
  \node[lbl,anchor=west] at (4.85853,0.95504) {$M_{5}$};
  \node[lbl,anchor=west] at (4.85853,1.19423) {$M_{4}$};
  \node[lbl,anchor=west] at (4.85853,1.7394) {$M_{6}$};
  \node[lbl,anchor=west] at (4.85853,1.97859) {$Y_0$};
  \node[lbl,anchor=west] at (4.85853,2.66624) {$M_{1}$};
  \node[lbl,anchor=west] at (4.85853,2.90543) {$M_{3}$};
  \node[lbl,anchor=west] at (4.85853,3.18952) {$M_{2}$};
  \node[lbl,anchor=west] at (4.85853,3.42871) {$L_{4}$};
  \node[lbl,anchor=west] at (4.85853,3.6679) {$R$};
  \node[lbl,anchor=south] at (-1.8107,4.13463) {$L_{5}$};
  \node[lbl,anchor=south] at (-1.42903,4.13463) {$L_{2}$};
  \node[lbl,anchor=south] at (0.02671,4.13463) {$L_{3}$};
  \node[lbl,anchor=south] at (0.36307,4.13463) {$L_{6}$};
  \node[lbl,anchor=north] at (-2.94889,-4.13463) {$R$};
  \node[lbl,anchor=north] at (-2.61253,-4.13463) {$L_{6}$};
  \node[lbl,anchor=north] at (-0.80117,-4.13463) {$L_{3}$};
  \node[lbl,anchor=north] at (-0.46481,-4.13463) {$L_{2}$};
  \node[lbl,anchor=north] at (2.30961,-4.13463) {$L_{5}$};
  \node[lbl,anchor=east] at (-4.85853,-0.45146) {$L_{4}$};
  \node[lbl,anchor=east] at (-4.85853,1.17519) {$M_{2}$};
  \node[lbl,anchor=east] at (-4.85853,1.41438) {$M_{3}$};
  \node[lbl,anchor=east] at (-4.85853,1.79005) {$M_{1}$};
  \node[lbl,anchor=east] at (-4.85853,2.02924) {$Y_0$};
  \node[lbl,anchor=east] at (-4.85853,2.47692) {$M_{6}$};
  \node[lbl,anchor=east] at (-4.85853,2.71611) {$M_{4}$};
  \node[lbl,anchor=east] at (-4.85853,2.97699) {$M_{5}$};
  \node[lbl,anchor=east] at (-4.85853,3.21618) {$L_{1}$};
\end{tikzpicture}
}
\caption{The starting arrangement after the doubling step \eqref{eq:new-lines}, with 54 triangular faces (8 of them $R$-triangles); the triple points are preserved.
The slopes $\mu_j$ are larger in absolute value than in \eqref{eq:new-lines}, since with the actual values the strip \eqref{eq:strip} would be indistinguishable at this scale, but the order of intersection points on each line is the same.}
\label{fig:double}
\end{figure}

The lines $L_1,\ldots,L_{n-2}$ are numbered in increasing order of their $x$-intercepts, so that $a_1<\cdots<a_{n-2}$. The lines $L_{n-1},L_n$ form the \emph{special pair}.
The doubling step imposes the following conditions on $\mathcal A$:
\begin{enumerate}[label=\textup{(D\arabic*)}]
\item the $x$-intercepts form the grid
\begin{equation}\label{eq:x-grid}
\begin{gathered}
\{a_1,\ldots,a_{n-2}\}
=\Bigl\{\tan\frac{k\pi}{n}:1\le |k|\le\frac n2-1\Bigr\},
\\
-\frac1n<a_{n-1}<0<a_n<\frac1n,
\end{gathered}
\end{equation}
in particular, $a_{n-2}=\cot(\pi/n)$;
\item the slopes of the special pair satisfy $m_{n-1}<0<m_n$;
\item each of the $n-1$ bounded segments of $Y_0$ is a side of exactly one triangular face of $\mathcal A$, and the faces on consecutive segments lie on opposite sides of $Y_0$.
\end{enumerate}
Under (D1) and (D2), the point $L_{n-1}\cap L_n$ lies below $Y_0$, and its $x$-coordinate lies in the interval $(a_{n-1},a_n)$.

\begin{lemma}[doubling step]\label{lem:doubling}
Let $n\ge2$ be even, and let $\mathcal A$ be a not necessarily simple arrangement satisfying \textup{(D1)--(D3)}. Then
\[
\mathcal B=\mathcal A\cup\{M_1,\ldots,M_n\}
\]
preserves all multiple points of $\mathcal A$ together with their multiplicities, has no new multiple points, and satisfies
\[
a_3(\mathcal B)=a_3(\mathcal A)+n^2.
\]
The arrangement $\mathcal B$ satisfies \textup{(D3)} with $2n$ in place of $n$. If, in addition,
\begin{equation}\label{eq:special-slopes}
|a_{n-1}|<\frac1{2n},\quad |a_n|<\frac1{2n},
\end{equation}
then $\mathcal B$ (after relabeling its lines) satisfies \textup{(D1)--(D3)} with $2n$ in place of $n$.
Finally,
\begin{equation}\label{eq:m-min}
m_{\min}(\mathcal B)<\frac{2\,m_{\min}(\mathcal A)}{n^{10}}<m_{\min}(\mathcal A).
\end{equation}
\end{lemma}

\begin{proof}
The computations in items~1--4 of the proof of~\cite[Prop.~3.1]{bbl} rely only on the equations of the lines: on condition \eqref{eq:x-grid}, on the signs of the slopes of the special pair,\footnote{However, item~4 of~\cite[Prop.~3.1]{bbl} needs a correction. There the signs of the slopes of the special pair are derived from the position of the point $L_{n-1}\cap L_n$ relative to $Y_0$. But the position of this point does not determine them: the two lines can also meet on that same side of $Y_0$ when their slopes have the same sign (then the $x$-coordinate of their intersection point lies outside $(a_{n-1},a_n)$). Therefore the signs $m_{n-1}<0<m_n$ are included in the hypotheses of the lemma as (D2), and everything the proof of~\cite[Prop.~3.1]{bbl} derives from the signs, namely $x(M_j\cap L_{n-1})<0<x(M_j\cap L_n)$, follows from the inequalities \eqref{eq:special-x} proved below.}
and on formulas \eqref{eq:new-lines}.
These computations do not use the simplicity of the arrangement. Therefore the estimates~\cite[(3)--(9)]{bbl}, the case analysis 3a--3c, and the resulting strict order of the intersection points on each line $M_j$ remain valid for our $\mathcal A$.
We restate the results of these computations as the following inequalities.
First, by~\cite[(6)]{bbl},
\begin{equation}\label{eq:L-M-coord}
|y(L_i\cap M_j)|<\frac{3m_{\min}}{n^9},
\qquad
|x(L_i\cap M_j)-a_i|<\frac3{n^9}.
\end{equation}
Second, by~\cite[(5)]{bbl}, the $y$-coordinates of the old intersection points off $Y_0$ are bounded away from zero:
\begin{equation}\label{eq:L-L-coord}
|y(L_i\cap L_j)|>\frac{m_{\min}}n,
\qquad \{i,j\}\neq\{n-1,n\},
\end{equation}
where the special pair $L_{n-1},L_n$ is excluded from the estimate because the $y$-coordinate of $L_{n-1}\cap L_n$ can be arbitrarily close to zero. Third, the case analysis 3a--3c and the estimates~\cite[(7)--(9)]{bbl} give the following:
\begin{equation}\label{eq:M-M-coord}
x(M_i\cap M_j)=
\begin{cases}
0, & \beta_i+\beta_j=0 \quad\text{(case 3a)},\\
-n^6\sin(2\beta_i), & \beta_i+\beta_j=\pm\pi/2 \quad\text{(case 3b)},\\
\tan(\beta_i+\beta_j)+\theta_{ij}, & \text{otherwise} \quad\text{(case 3c)},
\end{cases}
\end{equation}
where in case 3c
\begin{equation}\label{eq:theta}
\lvert\tan(\beta_i+\beta_j)\rvert<\frac n\pi,
\quad \frac3{n^8}<|\theta_{ij}|<\frac1{n^2},
\quad \operatorname{sgn}\theta_{ij}=\operatorname{sgn}\bigl((\beta_i+\beta_j)\beta_i\beta_j\bigr).
\end{equation}
Hence
\begin{equation}\label{eq:x-far}
|x(M_i\cap M_j)|\le n^6.
\end{equation}
Fourth,
\begin{equation}\label{eq:special-x}
x(M_j\cap L_{n-1})<a_{n-1}<0<a_n<x(M_j\cap L_n).
\end{equation}
Indeed,
\[
x(M_j\cap L_n)-a_n=\frac{\mu_j\,(a_n-b_j)}{m_n-\mu_j}.
\]
The denominator is positive: $|\mu_j|<2m_{\min}/n^{10}<m_{\min}\le m_n$ by~\cite[(3)]{bbl}.
Both factors $\mu_j$ and $a_n-b_j$ in the numerator have the sign of $-b_j$. For $\mu_j$ this is~\cite[(3)]{bbl} (for $\sigma=1$ the values $\mu_j$ and $b_j$ have the same sign, and here $\sigma=-1$). For $a_n-b_j$ it follows from $|b_j|\ge\tan\frac{\pi}{2n}>\frac1n>|a_n|$.
The computation for $L_{n-1}$ is similar.
Together with the signs of the slopes $m_{n-1}<0<m_n$, this implies that the new lines intersect the special pair strictly above $Y_0$.

Moreover, by~\cite[(3)]{bbl}, $|b_j|<2n/\pi$ and $|\mu_j|<2m_{\min}/n^{10}$.
The point $M_i\cap M_j$ lies on $M_j$, and its $y$-coordinate equals $\mu_j(x-b_j)$, where $x=x(M_i\cap M_j)$. For $n\ge2$ we have $|x-b_j|\le|x|+|b_j|<n^6+2n/\pi<2n^6$, and therefore
\[
|y(M_i\cap M_j)|=|\mu_j|\,|x-b_j|
<\frac{2m_{\min}}{n^{10}}\cdot2n^6=\frac{4m_{\min}}{n^4}.
\]
Since $4/n^4<1/n$ and $3/n^9<1/n$ for $n\ge2$, all intersection points involving at least one line $M_j$ lie in the strip
\begin{equation}\label{eq:strip}
|y|<\frac{m_{\min}}n.
\end{equation}
The threshold $m_{\min}/n$ is the same as in \eqref{eq:L-L-coord}. It separates the region where the old intersection points remain from the region where the new points appear.

Extending~\cite{bbl} to the non-simple case requires three more checks.

\emph{Old multiple points are preserved.}
Condition \eqref{eq:x-grid} gives $n$ distinct points $L_i\cap Y_0$, so all intersection points on $Y_0$ are simple, and every multiple point of $\mathcal A$ has the form $L_i\cap L_j$.
By \eqref{eq:L-L-coord}, it lies outside the strip \eqref{eq:strip}.
This holds for $L_{n-1}\cap L_n$ as well: this point lies below $Y_0$, so if it is multiple, it equals $L_k\cap L_{n-1}$ for some $k\notin\{n-1,n\}$, and \eqref{eq:L-L-coord} applies.
All intersection points on the new lines lie in the strip, so no new line passes through a multiple point of $\mathcal A$, and these points are preserved together with their multiplicities.

\emph{No new multiple points appear.}
A multiple point of $\mathcal B$ on a new line $M_j$ would mean that two intersection points on $M_j$ coincide.
But item~4 of the proof of~\cite[Prop.~3.1]{bbl} (with \eqref{eq:special-x} in place of the signs of the special pair) places the intersection points of $M_j$ with all other lines of $\mathcal B$ at pairwise distinct $x$-coordinates.
Also, $L_{n-1}\cap L_n$ is the only old intersection point off $Y_0$ that may lie in the strip \eqref{eq:strip}. Its $x$-coordinate lies in the interval $(a_{n-1},a_n)$, so by \eqref{eq:special-x} no new line passes through it.

\emph{Multiple points do not affect the triangle count.}
By \eqref{eq:L-L-coord}, the strip \eqref{eq:strip} contains no old intersection points except the $n$ intersection points on $Y_0$ and possibly $L_{n-1}\cap L_n$.
The new lines enter only the faces of $\mathcal A$ that have a side on $Y_0$.\footnote{In~\cite{bbl} this is stated with reference to estimates (5) and (6), which do not take the point $L_{n-1}\cap L_n$ into account. For $L\in\{L_{n-1},L_n\}$ this point lies below $Y_0$, while by \eqref{eq:special-x} $M_j$ meets $L$ above $Y_0$; for any other line $L$ through this point, the point is multiple and lies outside the strip (see above). Hence inside the strip every line of $\mathcal A$ meets $M_j$ on an edge ending on $Y_0$, so every face that $M_j$ enters has a side on $Y_0$.}
By \eqref{eq:L-M-coord}, the new lines intersect the old ones in the $3/n^{9}$-neighborhoods of the points $a_i=x(L_i\cap Y_0)$.
By (D3), the triangles with a side on $Y_0$ alternate between the two sides of $Y_0$, as they automatically do in a simple arrangement.
Therefore inside the strip the arrangement looks the same as in the simple case, and the count from the last paragraph of the same proof carries over unchanged: the new lines cut only the $n-1$ triangles with a side on $Y_0$, and each of them is replaced by $n+1$ triangles; another $n$ triangles appear at the two far ends of the strip; there are no other new triangles, and the remaining old ones are preserved.\footnote{
Strictly speaking, only $n-2$ of the triangles with a side on $Y_0$ disappear.
The triangle $Y_0L_{n-1}L_n$ bounded by the special pair is not replaced but remains a face of $\mathcal B$. It lies below $Y_0$, while by \eqref{eq:special-x} the new lines intersect $L_{n-1}$ and $L_n$ above $Y_0$ and intersect $Y_0$ outside the interval $[a_{n-1},a_n]$, since $|b_j|>1/n$.
Hence no new line crosses its boundary, and therefore no new line enters it.
However, $n$ new triangles are formed between the lines of the special pair, so this remark does not affect the triangle count.
Also, in the intuitive picture (\cref{fig:double}), where $Y_0$ ``thickens'' into the set $\{Y_0, M_1, \ldots, M_n\}$, this preserved triangle does not differ from the disappearing ones.}
Consequently,
\[
a_3(\mathcal B)
=a_3(\mathcal A)-(n-1)+(n-1)(n+1)+n
=a_3(\mathcal A)+n^2.
\]
The argument that each bounded segment of $Y_0$ is a side of at least one triangular face of $\mathcal B$ also carries over.
A segment cannot be a side of two such faces, since both its endpoints are intersection points of multiplicity 2 (see \cref{lem:double-segment}).
Suppose that the triangles on two consecutive segments of $Y_0$ lie on the same side of $Y_0$. Both have a side on the line $L$ through the common endpoint of the segments, and this side is the edge of $L$ that starts at $Y_0$ on that side. Hence the third vertex of both triangles is the other endpoint of this edge, and the three lines through the three consecutive intersection points on $Y_0$ meet there.
In $\mathcal B$ this is impossible: the points $b_j$ alternate on $Y_0$ with the points $a_1,\ldots,a_{n-2}$, and the pair $a_{n-1},a_n$ lies between $b_j=\pm\tan\frac\pi{2n}$, since $\tan\frac\pi{2n}>\frac1n$. So among any three consecutive points at least one lies on a new line, and the new lines have no multiple points.

Finally,
\[
\{a_1,\ldots,a_{n-2}\}\cup\{b_1,\ldots,b_n\}
=\left\{\tan\frac{k\pi}{2n}:1\le|k|\le n-1\right\}.
\]
If \eqref{eq:special-slopes} holds, then, after relabeling, $a_{n-1}$ and $a_n$ remain the $x$-intercepts of the special pair, and the others, numbered in increasing order, satisfy \eqref{eq:x-grid} with $2n$ in place of $n$.
Since the same two lines remain the special pair, condition (D2) and the sign $\sigma=-1$ are inherited.
The estimate \eqref{eq:m-min} follows from $|\mu_j|<2m_{\min}/n^{10}$, whose right-hand side is less than $m_{\min}$ for $n\ge2$. So the minimum $m_{\min}(\mathcal B)$ is attained at one of the new lines.
\end{proof}

\Cref{lem:doubling} describes the doubling without the line $R$, as in~\cite{bbl}.
When $R$ is restored, the new lines also change the $R$-triangles, and we now count these changes.
Since $R$ lies far to the right, it cuts off triangles from pairs of lines that meet at their last intersection point.

We call a pair of lines of an arrangement a \emph{right wedge} (\cref{fig:double-zoom}) if their common point (the \emph{apex}) is
the rightmost intersection point on each of them. Beyond this point, neither line
meets any other line of the arrangement.
By \eqref{eq:x-grid}, the rightmost intersection point of the old lines $L_i$ with $Y_0$ has $x$-coordinate $a_{n-2}=\cot(\pi/n)$.
For the count of $R$-triangles, the doubling step imposes two more conditions on $\mathcal A$:
\begin{enumerate}[label=\textup{(W\arabic*)}]
\item $(Y_0,L_{n-2})$ is a right wedge of $\mathcal A$;
\item no $R$-triangle of $\mathcal A\cup\{R\}$, except the triangle $\Delta_n$ bounded by $R$, $Y_0$, and $L_{n-2}$, meets the strip \eqref{eq:strip}.
\end{enumerate}

\begin{figure}[H]
\centering
\resizebox{\ifdim\width>\linewidth\linewidth\else\width\fi}{!}{\begin{tikzpicture}[x=76.8836cm,y=76.8836cm]
  \begin{scope}
    \clip (-0.08454,-0.01047) rectangle (0.08454,0.01047);
    \fill[rface] (0.05284,0.00403) -- (0.04479,-0.00403) -- (-0.05284,0.00403) -- cycle;
    \fill[face] (-2.05284,0.00403) -- (-0.05284,0.00403) -- (-1.71086,0.14082) -- cycle;
    \draw[rline] (0.03834,-0.01047) -- (0.05928,0.01047);
    \draw[oldline] (-0.08454,0.00403) -- (0.08454,0.00403);
    \draw[newline] (-0.08454,0.00664) -- (0.08454,-0.00731);
  \end{scope}
  \node[lbl,anchor=west] at (0.08454,-0.00731) {$M_{6}$};
  \node[lbl,anchor=west] at (0.08454,0.00403) {$Y_0$};
  \node[lbl,anchor=south] at (0.05928,0.01047) {$R$};
\end{tikzpicture}
}
\caption{The $R$-triangle at the point $R\cap Y_0$ in \cref{fig:double} ($n=6$), enlarged. The line $R$ cuts it off from the right wedge $(Y_0,M_6)$. This is the triangle $\Delta_{12}$ of \textup{(W2)} at the next step, with $M_6$ in the role of $L_{10}$.}
\label{fig:double-zoom}
\end{figure}

The following lemma lists the right wedges of $\mathcal B$ that involve a new line, for $n\equiv0\pmod4$.
In our series, this case covers all iterations with $s\ge1$: $n_s=6\cdot2^s\equiv0\pmod{12}$.
The only remaining case $s=0$ (that is, $n=6$) is treated in \cref{lem:first-step}.

\begin{lemma}\label{lem:right-wedges}
Let $n\ge8$, with $n\equiv0\pmod4$, let the arrangement $\mathcal A$ satisfy \textup{(D1)--(D3)} and \textup{(W1)}, let the slope of the line $L_{n-2}$ be negative, and let $\mathcal B=\mathcal A\cup\{M_1,\ldots,M_n\}$.
Then the right wedges of $\mathcal B$ that involve at least one new line are exactly the pairs
\begin{gather}
(M_i, M_{n/2+1-i}),
\qquad 1\le i\le\frac n4,\label{eq:wedges-low}\\
(M_i,M_{3n/2-i}),
\qquad \frac n2<i<\frac{3n}{4},\label{eq:wedges-high}\\
(L_{n-2},M_{3n/4}),
\qquad
(Y_0,M_n).\label{eq:wedges-old}
\end{gather}
There are $n/2+1$ of them, and each new line belongs to exactly one.
Moreover, $\mathcal B$ satisfies \textup{(W1)} with $2n$ in place of $n$ and $M_n$ in the role of $L_{2n-2}$.
\end{lemma}

\begin{proof}
Among all pairs of new lines $(M_i,M_j)$, $i\neq j$, consider those with $\beta_i+\beta_j=\pm\pi/2$.
By \eqref{eq:new-lines}, $\beta_i+\beta_j=-\pi+(i+j-1)\pi/n$, so
\begin{gather}
\beta_i+\beta_j=-\pi/2 \iff i+j=n/2+1,\label{eq:Par1}\\
\beta_i+\beta_j=\pi/2 \iff i+j=3n/2+1.\label{eq:Par2}
\end{gather}
Each of these equalities implies $b_ib_j=1$ and $\sin(2\beta_i)=\sin(2\beta_j)$.
For such values the slopes $\mu_i$ and $\mu_j$ differ only in the term with $n^{-6}$. Hence the angle between the lines $M_i,M_j$ is small, and they intersect far from the other intersection points of the arrangement.
We call such a pair \emph{almost parallel}. It corresponds to case 3b of~\cite[Prop.~3.1]{bbl}, and by \eqref{eq:M-M-coord} $x(M_i\cap M_j)=-n^6\sin(2\beta_i)$.

For a pair of lines satisfying \eqref{eq:Par1}, both $\beta_i$ and $\beta_j$ are negative, so $\sin(2\beta_i)<0$ and $x(M_i\cap M_j)>0$. Together with the inequality $\lvert\sin(2\beta_i)\rvert>\pi/2n$ from~\cite[(3)]{bbl}, this gives $x(M_i\cap M_j)>\pi n^5/2>n^5$.
For a pair of lines satisfying \eqref{eq:Par2}, both $\beta_i$ and $\beta_j$ are positive, and a similar estimate gives $x(M_i\cap M_j)<-n^5$, so these intersection points lie far to the left and are not the rightmost points on their lines.

Let $\beta_i<0$, that is, $i\le n/2$.
Then $M_i$ has an almost parallel partner $M_j$ with $\beta_j$ determined by \eqref{eq:Par1}: $j=n/2+1-i$, and $j\neq i$, since $n/2+1$ is odd for $n\equiv0\pmod4$.
All other intersection points of $M_i$ lie in the region $|x|<n$: $x(M_i\cap Y_0)=b_i$, $x(M_i\cap L_k)$
is bounded by \eqref{eq:L-M-coord}, and the bounds on the intersection points of $M_i$ with the other new lines are given by cases 3a and 3c of~\cite[Prop.~3.1]{bbl}.
Therefore the intersection point $M_i\cap M_j$ with $x$-coordinate greater than $n^5$ is the rightmost on $M_i$.
The angle $\beta_j$ is also negative, and the same argument applies to $M_j$: its partner from \eqref{eq:Par1} is $M_i$.
Hence the pairs \eqref{eq:wedges-low} form right wedges.
With \eqref{eq:Par1}, the condition $i<j$, which orders the lines in a pair and lists each pair once, reduces to $i\le n/4$.

Now let $\beta_i>0$, that is, $i>n/2$.
As shown above, for the almost parallel line $M_j$ from \eqref{eq:Par2} we have $x(M_i\cap M_j)<-n^5$, so their intersection point lies far to the left.
We look for the rightmost of the remaining intersection points on $M_i$. In case 3c of~\cite[Prop.~3.1]{bbl}, by \eqref{eq:M-M-coord} and \eqref{eq:theta}, the intersection with a new line $M_j$ has $x=\tan(\beta_i+\beta_j)+\theta_j$ with $3/n^8<|\theta_j|<1/n^2$. In case 3a, it has exactly $x=0=\tan(\beta_i+\beta_j)$. The intersection with an old line $L_k$ lies in the $3/n^{9}$-neighborhood of the point $a_k$ by \eqref{eq:L-M-coord}. The intersection with $Y_0$ is exactly the point $b_i$.

Since $\beta_i>0$, the sums $z=\beta_i+\beta_j$ lie in the interval $(-\pi/2,\pi)$.
For $z\in(\pi/2,\pi)$ the tangent is negative, and such points lie to the left of zero and hence to the left of $a_{n-2}$.
On the interval $(-\pi/2,\pi/2)$ the tangent increases, so the rightmost point corresponds to the largest sum in this interval.
On the grid \eqref{eq:new-lines} the sums are spaced $\pi/n$ apart, and the value $z=\pi/2$ corresponds to \eqref{eq:Par2}, that is, to the almost parallel partner, whose intersection point lies far to the left.
Therefore the largest available sum is $\pi/2-\pi/n$.
It corresponds to the partner $M_{3n/2-i}$ (if $3n/2-i\neq i$) and gives
\[
x(M_i\cap M_{3n/2-i})=\tan\left(\frac\pi2-\frac\pi n\right)+\theta
=a_{n-2}+\theta,
\qquad \frac3{n^8}<|\theta|<\frac1{n^2} .
\]
Here we used $\tan(\pi/2-\pi/n)=\cot(\pi/n)=a_{n-2}$. By \eqref{eq:theta}, the sign of $\theta$ is the sign of the product $(\beta_i+\beta_j)\beta_i\beta_j$. For $i\le n-1$ all three factors are positive, since we consider $\beta_i>0$, and $\beta_j=\pi/2-\pi/n-\beta_i>0$; for $i=n$ we have $\beta_j<0$.
Among the intersections of $M_i$ with the new lines, the rightmost neighborhood is that of the point $a_{n-2}$. By \eqref{eq:L-M-coord}, the intersection of $M_i$ with the rightmost old line $L_{n-2}$ falls into the neighborhood of the same point (for the other old lines $L_k$ we have $a_k<a_{n-2}$).
Within this neighborhood, for $i\le n-1$ the intersection with $M_{3n/2-i}$ lies farther to the right, since its offset $\theta$ is greater than $3/n^{8}$, while the offset of the point $M_i\cap L_{n-2}$ is less than $3/n^{9}$ in absolute value.
It remains to take the line $Y_0$ into account. Its intersection point with $M_i$ has $x$-coordinate $b_i$ and lies in none of the neighborhoods listed, so it has to be compared with $a_{n-2}$ separately.
The values $b_i$ increase with $i$.
For $i<n$ we have
\[
a_{n-2}-b_i\ge a_{n-2}-b_{n-1}=\cot\frac\pi n-\cot\frac{3\pi}{2n}=\frac{\sin(\pi/2n)}{\sin(\pi/n)\sin(3\pi/2n)}\ge\sin\frac\pi{2n}\ge\frac1n,
\]
so the point $M_i\cap Y_0$ lies to the left of the neighborhood of $a_{n-2}$ and is not the rightmost point on $M_i$.
For $i=n$, by the identity $\cot(x/2)=\cot x+1/\sin x$,
\begin{equation}\label{eq:b-n}
b_n=\cot\frac\pi{2n}=a_{n-2}+\frac1{\sin(\pi/n)}>a_{n-2}+\frac n\pi,
\end{equation}
so on the line $M_n$ the point $M_n\cap Y_0$ is the rightmost.
It is the rightmost on $Y_0$ as well: the other intersection points on $Y_0$ have $x$-coordinates $b_j<b_n$ and $a_k\le a_{n-2}<b_n$.
Hence $(Y_0,M_n)$ is a right wedge, the second pair in \eqref{eq:wedges-old}.

The argument with the partner $M_{3n/2-i}$ has one exception, $i=3n/4$, when $3n/2-i=i$ and there is no partner with the sum $\pi/2-\pi/n$.
The next available sum $\pi/2-2\pi/n$ gives only $x\le\cot(2\pi/n)+1/n^{2}$.
This is to the left of the neighborhood of $a_{n-2}$, which has radius less than $1/n^{2}$, since
\[
a_{n-2}-\cot\frac{2\pi}n=\frac1{\sin(2\pi/n)}\ge1 .
\]
Therefore on the line $M_{3n/4}$ the rightmost intersection point is the one with the old line $L_{n-2}$.

Thus the rightmost intersection point on $M_i$ is the one with the line $M_{3n/2-i}$ for $i\notin\{3n/4,n\}$, with $L_{n-2}$ for $i=3n/4$, and with $Y_0$ for $i=n$.
Accordingly, the right wedges are the pairs \eqref{eq:wedges-high} and the two pairs \eqref{eq:wedges-old}.
For the pairs \eqref{eq:wedges-high}, the argument above suffices: both lines have indices greater than $n/2$, and the common point of the pair is the rightmost on each of them.
The pair $(Y_0,M_n)$ has also been checked on both lines.

We check that the common point of the pair $(L_{n-2},M_{3n/4})$ is the rightmost on $L_{n-2}$ as well.
By the hypothesis of the lemma,
\[
L_{n-2}:y=\ell(x-a_{n-2}),
\qquad \ell<0.
\]
The new lines are almost horizontal compared with $L_{n-2}$: $|\mu_j|<2m_{\min}/n^{10}<|\ell|$ by~\cite[(3)]{bbl}, since the minimum in $m_{\min}(\mathcal A)$ is also taken over $\ell$.
Hence $L_{n-2}$ lies above each line $M_j$ to the left of their common point and below it to the right.
Let $U=L_{n-2}\cap M_{3n/4}$. The line $L_{n-2}$ meets $M_k$ to the right of $U$ only if $M_k$ passes below $M_{3n/4}$ at $x=x(U)$.
So it suffices to show that at $x=x(U)$ the line $M_{3n/4}$ is the lowest of the new lines.

By \eqref{eq:new-lines}, the line $M_j$ passes at height
\begin{equation}\label{eq:M-height}
\mu_j(a_{n-2}-b_j)=-\frac{m_{\min}}{n^{10}}F_j,
\quad\!\!
F_j=\Bigl(\sin(2\beta_j)+\frac1{n^6\tan\beta_j}\Bigr)\left(a_{n-2}-\tan\beta_j\right),
\end{equation}
at $x=a_{n-2}$.
By \eqref{eq:L-M-coord}, $|x(U)-a_{n-2}|<3/n^9$, so between $a_{n-2}$ and $x(U)$ the height of each new line changes by less than $|\mu_j|\cdot3/n^9<6m_{\min}/n^{19}$.
Therefore $M_{3n/4}$ is the lowest at $x=x(U)$ if at $x=a_{n-2}$ it lies below every other new line by more than $12m_{\min}/n^{19}$, that is, by \eqref{eq:M-height}, if
\begin{equation}\label{eq:F-gap}
F_{3n/4}-F_k>\frac{12}{n^9}
\qquad\text{for all }k\neq\frac{3n}4.
\end{equation}

We split $F_j=G_j+H_j$ according to the two terms in parentheses:
\[
G_j=\sin(2\beta_j)\left(a_{n-2}-\tan\beta_j\right),
\qquad
H_j=\frac{a_{n-2}-\tan\beta_j}{n^{6}\tan\beta_j}.
\]
For $G_j$ we use the identity
\begin{equation}\label{eq:tan-id}
\sin(2\beta)(a_{n-2}-\tan\beta)
=\frac{\sin(2\beta+\pi/n)}{\sin(\pi/n)}-1 .
\end{equation}
By \eqref{eq:new-lines}, $2\beta_j+\pi/n=2\pi j/n-\pi$, so
\begin{equation}\label{eq:G-j}
G_j=-\frac{\sin(2\pi j/n)}{\sin(\pi/n)}-1 ,
\end{equation}
and the maximum over $j$ corresponds to the minimum of $\sin(2\pi j/n)$, that is, to $2\pi j/n=3\pi/2$:
\[
j=\frac{3n}{4},
\qquad
\beta_j=\frac\pi4-\frac\pi{2n},
\]
and the maximum equals $1/\sin(\pi/n)-1$.
For $j=3n/4\pm1$ we have $\sin(2\pi j/n)=\sin(3\pi/2\pm2\pi/n)=-\cos(2\pi/n)$, so
the maximum exceeds the values at these indices by
\[
\frac1{\sin(\pi/n)}-1-\frac{\cos(2\pi/n)}{\sin(\pi/n)}+1=
\frac{1-\cos(2\pi/n)}{\sin(\pi/n)}
=\frac{2\sin^2(\pi/n)}{\sin(\pi/n)}
=2\sin\frac\pi n\ \ge\ \frac4n .
\]
Since $|a_{n-2}-\tan\beta_j|<3n/\pi<n$~\cite[Prop.~3.1, item~2]{bbl} and
$\lvert\tan\beta_j\rvert=|b_j|>\pi/2n$ by~\cite[(3)]{bbl},
\[
|H_j|=\frac{|a_{n-2}-\tan\beta_j|}{n^6\,|\tan\beta_j|}<\frac n{n^6}\cdot\frac{2n}\pi<\frac1{n^4}.
\]
Therefore, for $k\ne3n/4$,
\[
F_{3n/4}-F_k>\frac4n-\frac2{n^4}>\frac{12}{n^9},
\]
which proves \eqref{eq:F-gap}. Hence $U$ is the rightmost point on $L_{n-2}$ among the intersection points with the new lines.

The old intersection points on $L_{n-2}$ lie to the left of $U$.
Indeed, $F_{3n/4}>0$, since both factors in \eqref{eq:M-height} are positive for $\beta_{3n/4}=\pi/4-\pi/2n$. So $M_{3n/4}$ passes below the point $(a_{n-2},0)$ of $L_{n-2}$, and $U$ lies to the right of $a_{n-2}$. By (W1), there are no old intersection points of $\mathcal A$ on $L_{n-2}$ to the right of $a_{n-2}$.
Hence $U$ is the rightmost intersection point of $\mathcal B$ on $L_{n-2}$, and the pair $(L_{n-2},M_{3n/4})$ is indeed a right wedge.

Thus the rightmost intersection point is found for each new line, and the indices in \eqref{eq:wedges-low}--\eqref{eq:wedges-old} cover $1,\ldots,n$ exactly once: \eqref{eq:wedges-low} gives $1,\ldots,n/2$, \eqref{eq:wedges-high} gives $n/2+1,\ldots,3n/4-1$ and $3n/4+1,\ldots,n-1$, and \eqref{eq:wedges-old} gives the remaining $3n/4$ and $n$.
By \cref{lem:doubling}, the new lines have no multiple points. Hence each of them has a single rightmost intersection point and belongs to at most one wedge.
Therefore \eqref{eq:wedges-low}--\eqref{eq:wedges-old} are all the right wedges that involve at least one new line, and there are exactly $n/2+1$ of them.

By \eqref{eq:b-n}, in \textup{(W1)} the role of $L_{n-2}$ with $a_{n-2}=\cot(\pi/n)$ passes to $M_n$ with $b_n=\cot(\pi/2n)$ (\cref{fig:double-zoom}).
\end{proof}

\begin{lemma}\label{lem:r-triangle}
Let $n\ge6$, with either $n\equiv0\pmod4$ or $n=6$, let the arrangement $\mathcal A$ satisfy \textup{(D1)--(D3)}, \textup{(W1)}, \textup{(W2)}, and let $\mathcal B=\mathcal A\cup\{M_1,\ldots,M_n\}$.
Suppose also that the slope of the line $L_{n-2}$ is at most $1/2$ in absolute value, negative for $n\equiv0\pmod4$ and positive for $n=6$, and that the line $R$ has the form $y=x-\delta$ with $\delta\ge2n^6$.
Then $\Delta_n$ is a face of $\mathcal A\cup\{R\}$, and in $\mathcal B\cup\{R\}$ it is replaced by exactly $n/2+1$ new $R$-triangles, while all other $R$-triangles of $\mathcal A\cup\{R\}$ remain faces.
Consequently,
\[
r(\mathcal B\cup\{R\})
=r(\mathcal A\cup\{R\})+\frac n2.
\]
Moreover, $\mathcal B$ satisfies \textup{(W2)} with $2n$ in place of $n$ and $M_n$ in the role of $L_{2n-2}$.
\end{lemma}

\begin{proof}
By \cref{lem:right-wedges} for $n\equiv0\pmod4$ and by \cref{lem:first-step} for $n=6$, the right wedges of $\mathcal B$ that involve at least one new line are the $n/2+1$ pairs \eqref{eq:wedges-low}--\eqref{eq:wedges-old} and \eqref{eq:wedges-first-step}, respectively, and each new line belongs to exactly one of them.
Other right wedges, formed only by old lines $L_i$, give no new $R$-triangles, and we do not consider them.

\emph{Step 1. The new $R$-triangles are exactly the triangles that $R$ cuts off from these wedges.} For each line $D$ that belongs to a right wedge,
$x(R\cap D)>n^6$.
Indeed, by hypothesis $R:y=x-\delta$ with $\delta\ge 2n^6$, and the slopes of all lines in the wedges are at most $1/2$ in absolute value: $Y_0$ has slope zero, the slope of $L_{n-2}$ is bounded by the hypothesis of the lemma, and the new lines have $|\mu_j|<2m_{\min}(\mathcal A)/n^{10}\le 1/n^{10}$ by~\cite[(3)]{bbl}, since $m_{\min}(\mathcal A)\le|m_{n-2}|\le1/2$.
The $x$-intercepts $a$ in the equations of these lines are less than $n$ in absolute value (they are $b_j$ or $a_{n-2}$), so
\[
x(R\cap D)=\frac{\delta-ma}{1-m}\ge\frac{\delta-n/2}{3/2}\ge\frac43n^6-\frac n3>n^6 .
\]

The apexes of the wedges in $\mathcal{B}$ have $x$-coordinates at most $n^6$. For a pair of new lines, the apex is a point $M_i\cap M_j$, which satisfies \eqref{eq:x-far}. For the two pairs with an old line, the apex has $x$-coordinate less than $n$.
For the pair with the line $L_{n-2}$, the apex is a point $L_{n-2}\cap M_j$, and by \eqref{eq:L-M-coord} and~\cite[(4)]{bbl}
\[
x(L_{n-2}\cap M_j)<a_{n-2}+\frac3{n^9}=\cot\frac\pi n+\frac3{n^9}<\frac n\pi+1 .
\]
For the pair $(Y_0,M_n)$, the apex is the point $Y_0\cap M_n$, and by~\cite[(3)]{bbl}
\[
x(Y_0\cap M_n)=b_n=\cot\frac\pi{2n}<\frac{2n}\pi .
\]
Both expressions are less than $n$ for $n\ge6$. Therefore $R$ intersects the lines of the wedge $(D,D')$ on the rays to the right of the apex, and the three lines bound a nondegenerate triangle $T$.

This triangle is a face of the arrangement $\mathcal{B}\cup \{R\}$.
Indeed, the rays forming the wedge extend to the right from its apex $Z=D\cap D'$ without crossing the other lines.
Therefore a line $D''$ of the arrangement can cross the sides of $T$ on $D$ and $D'$ only at the apex $Z$ itself.
But only two lines pass through $Z$, since each of these wedges involves at least one new line $M_j$, and by \cref{lem:doubling} the new lines have no multiple points.
Hence $D''$ crosses the boundary of $T$ at most once (possibly on the line $R$) and cannot cut off any part of $T$.

Thus $R$ closes each of the $n/2+1$ right wedges into a triangular face of $\mathcal B\cup\{R\}$.
The same argument applies to the wedge $(Y_0,L_{n-2})$ of $\mathcal A$: its apex is simple, since all intersection points on $Y_0$ are distinct by \eqref{eq:x-grid}, and the slopes $0$ and $m_{n-2}$ are at most $1/2$ in absolute value.
Hence $\Delta_n$ is also a face of $\mathcal A\cup\{R\}$, which proves the first claim of the lemma.

There are no other new $R$-triangles. Indeed, a face of $\mathcal B\cup\{R\}$ bounded by three lines of $\mathcal A\cup\{R\}$ is also a face of $\mathcal A\cup\{R\}$, since adding lines only subdivides faces. Hence every new $R$-triangle has a side on some new line $M_j$.
This side goes from the intersection point with some line $D$ to the point $M_j\cap R$, which lies to the right of all other intersection points on $M_j$. The side has no intersection points in its interior only if $M_j\cap D$ is the rightmost intersection point on $M_j$.
This point is simple by \cref{lem:doubling} and is the apex of the wedge containing $M_j$, so $(D,M_j)$ is one of these wedges.
Hence there are exactly $n/2+1$ new $R$-triangles. They replace the single face $\Delta_n$ (its side on $Y_0$ is crossed by $M_n$ at the point $b_n>a_{n-2}$), which gives an increase of $n/2$.

\emph{Step 2. The old $R$-triangles other than $\Delta_n$ remain faces.} Let $F$ be an $R$-triangle of $\mathcal A\cup\{R\}$ other than $\Delta_n$, and suppose that a new line $M_j$ enters the interior of $F$.
The face $F$ is convex and bounded, so $M_j$ crosses its boundary at exactly two points.
The side of $F$ on $R$ contains at most one of them, since $M_j$ and $R$ intersect once.
Hence at least one of these two points lies on a side belonging to some line $L$ of $\mathcal A$. This point is the intersection point $M_j\cap L$, and by \eqref{eq:strip} it lies in the strip $|y|<m_{\min}(\mathcal A)/n$.
But this point belongs to $F$, and by \textup{(W2)} $F$ does not meet this strip, a contradiction.
Thus no new line enters the interior of $F$, and $F$ remains a face of $\mathcal B\cup\{R\}$.

\emph{Step 3. The new $R$-triangles other than $\Delta_{2n}$ avoid the next strip, so \textup{(W2)} is inherited.} Let $m'=m_{\min}(\mathcal B)$, and let $S$ be the strip $|y|<m'/2n$, that is, the strip \eqref{eq:strip} for $\mathcal B$ at the next step.
For the old $R$-triangles the required condition already holds: by \eqref{eq:m-min}, $m'<m_{\min}(\mathcal A)$, so $S$ is contained in the strip $|y|<m_{\min}(\mathcal A)/n$, which they do not meet by \textup{(W2)}.

Now let $T$ be a new $R$-triangle corresponding to a wedge $(D,D')$ other than $(Y_0,M_n)$. Its vertices are $Z=D\cap D'$, $R\cap D$, and $R\cap D'$.
Neither $D$ nor $D'$ coincides with $Y_0$, so the point $D\cap Y_0=(a_D,0)$ is one of the intersection points on $D$, and $Z$ is the rightmost of them. The vertex $Z$ is simple, so $Z\neq D\cap Y_0$ and $x(Z)>a_D$.
Similarly, $x(Z)>a_{D'}$.
Therefore both factors in parentheses in the equality
\[
y(Z)=m_D\,\bigl(x(Z)-a_D\bigr)=m_{D'}\bigl(x(Z)-a_{D'}\bigr)
\]
are positive, so the slopes $m_D$ and $m_{D'}$ have the same sign.
Denote this sign by $\tau$, so that $\tau\,y(Z)>0$.
Further, $x(R\cap D)>n^6>n>|a_D|$ by the estimate of Step 1, so $\tau\,y(R\cap D)>0$, and the same holds for $D'$.
Thus all three vertices of $T$ lie on the same side of $Y_0$.

We estimate their distance from $Y_0$.
The slopes of both lines of the wedge are at least $m'$ in absolute value (both lines differ from $Y_0$) and at most $1/2$ by the estimate of Step 1, so $1-m_D\in[1/2,3/2]$ and
\[
\bigl|y(R\cap D)\bigr|=\frac{|m_D|\,(\delta-a_D)}{1-m_D}
\ \ge\ \frac{2\,m'\,(2n^6-n)}3\ >\ m'n^6 ,
\]
and the same holds for $D'$.
For the vertex $Z$, we take the line of the wedge that is new: each of the $n/2+1$ wedges except $(Y_0,M_n)$ involves a line $M_j$ with $j\le n-1$.
The point $M_j\cap L_{n-2}$ lies on $M_j$, and $Z$ is the rightmost intersection point on $M_j$, so by \eqref{eq:L-M-coord}
\[
x(Z)\ \ge\ x(M_j\cap L_{n-2})\ >\ a_{n-2}-\frac3{n^9} .
\]
The values $b_j=\tan\beta_j$ increase with $j$, so $b_j\le b_{n-1}=\cot\frac{3\pi}{2n}$ and
\[
x(Z)-b_j\ >\ \cot\frac\pi n-\cot\frac{3\pi}{2n}-\frac3{n^9}
=\frac{\sin\frac\pi{2n}}{\sin\frac\pi n\,\sin\frac{3\pi}{2n}}-\frac3{n^9}
\ >\ \frac{2n}{3\pi^2}-\frac3{n^9}\ >\ \frac n{16} ,
\]
where we used $\sin x>2x/\pi$ on $(0,\pi/2)$ for the numerator and $\sin x<x$ for the denominator.
Hence
\[
|y(Z)|=|\mu_j|\,\bigl(x(Z)-b_j\bigr)>\frac{m'n}{16} .
\]
Both bounds, $m'n^6$ and $m'n/16$, exceed $m'/2n$ for $n\ge6$, so all three vertices of $T$ lie in the half-plane $\tau y\ge m'/2n$. This half-plane is convex, so the whole triangle $T$ lies in it and does not meet the strip $S$.

The triangle $\Delta_{2n}=RY_0M_n$, by contrast, meets the strip $S$: two of its vertices, $R\cap Y_0=(\delta,0)$ and $Y_0\cap M_n=(b_n,0)$, lie on the line $Y_0$ itself.
\end{proof}

We need to count the triangular faces of the arrangement $\mathcal A_s\cup\{R\}$.
\Cref{lem:r-avoids} below shows that the line $R$ cuts no triangular face of $\mathcal A_s$.
Then the numbers of triangular faces of $\mathcal A_s$ and of $R$-triangles can simply be added.

\begin{lemma}\label{lem:r-avoids}
For each $0\le s\le t$, the line $R$ passes through no intersection point of $\mathcal A_s$ and has no common points with the interior of any triangular face of $\mathcal A_s$.
Consequently, adding $R$ preserves the multiplicities of all points, all triangular faces of $\mathcal A_s$ remain faces of $\mathcal A_s\cup\{R\}$, and
\begin{equation}\label{eq:a3-with-R}
a_3(\mathcal A_s\cup\{R\})
=a_3(\mathcal A_s)+r(\mathcal A_s\cup\{R\}).
\end{equation}
\end{lemma}

\begin{proof}
Denote by $V=L_2\cap L_3$ the intersection point of the starting arrangement on the two lines with large slopes $\mp\eta$. The point $V$ is simple in $\mathcal A_0$ by \cref{lem:seed}, and hence in all $\mathcal A_s$ by \cref{lem:doubling}. Set
\[
\Pi=\left\{(x,y):|x|\le\frac{n_t^6}{64},\ |y|\le\frac{n_t^6}{64}\right\}.
\]

\emph{All intersection points of $\mathcal A_s$ except $V$ lie in $\Pi$.}
For the intersection points of the starting arrangement, this is \cref{lem:seed-box}: all points except $V$ have both coordinates less than 14 in absolute value, and $n_t^6/64\ge6^6/64=729$.
An intersection point added at the step with $n=n_{s'}$, $s'<s$, lies on at least one new line $M_j$. Therefore by \eqref{eq:strip} its $y$-coordinate is less than $m_{\min}(\mathcal A_{s'})/n_{s'}<1$ in absolute value, and its $x$-coordinate does not exceed $n_{s'}^6\le(n_t/2)^6=n_t^6/64$ in absolute value.
Indeed, for points of the form $M_i\cap M_j$, the $x$-coordinate is bounded by \eqref{eq:x-far}. The point $M_j\cap Y_0$ has $x$-coordinate $b_j$ with $|b_j|<2n_{s'}/\pi$ by~\cite[(3)]{bbl}, and the $x$-coordinate of the point $M_j\cap L_i$ differs from $a_i$ by less than $3/n_{s'}^9$ by \eqref{eq:L-M-coord}, where $|a_i|\le\cot(\pi/n_{s'})<n_{s'}/\pi$ by \eqref{eq:x-grid}. So in the last two cases the $x$-coordinate is less than $n_{s'}$ in absolute value.

\emph{$R$ does not enter $\Pi$ and contains no intersection points of $\mathcal A_s$.}
In $\Pi$ we have
\[
x-y\le \frac{n_t^6}{32}<2n_t^6=\delta,
\]
whereas $x-y=\delta$ everywhere on $R$. For the point $V$, the difference $x-y$ equals $3n_t^6\neq\delta$. Hence $R$ contains no intersection points of $\mathcal A_s$: adding $R$ preserves the multiplicities, and all points on $R$ in $\mathcal A_s\cup\{R\}$ are simple.

\emph{Faces without the vertex $V$.} A bounded face of a line arrangement is a convex polygon, hence the convex hull of its vertices, and each of its vertices is an intersection point of lines of the arrangement.
The set $\Pi$ is also convex. Hence a face all of whose vertices lie in $\Pi$ lies entirely in $\Pi$ and does not meet $R$.

\emph{No triangular face has the vertex $V$, for any $s$.}
Here the lines $L_1,\ldots,L_6$ keep their labels from the starting arrangement.
On the line $L_2$ the slope is $-\eta<0$, so $x=-1/\sqrt3-y/\eta$, and the order of the intersection points by $x$-coordinate is opposite to their order by $y$-coordinate.
By \cref{cor:rightmost-L2}, the point $V$ is the rightmost intersection point of $\mathcal A_0$ on $L_2$, and its neighbor is $L_2\cap L_6$. Hence $y(L_2\cap L_6)$ is the smallest $y$-coordinate among the intersection points of $\mathcal A_0$ on $L_2$ other than $V$, and $y(L_2\cap L_6)<-1$ by \cref{lem:seed-box}.
Every new intersection point on $L_2$, at whatever step it appears, lies in the strip \eqref{eq:strip} of its step.
These strips shrink: at step $s'$ the strip is $|y|<m_{\min}(\mathcal A_{s'})/n_{s'}$, where the denominator grows with $s'$ and the numerator decreases by \eqref{eq:m-min}.
The widest strip corresponds to the step $s'=0$, where $m_{\min}(\mathcal A_0)=2/5$ by \cref{lem:uv} and $n_0=6$, and gives $|y|<1/15$.
Therefore the new intersection points lie above $L_2\cap L_6$, that is, to the left of it, and the neighbor of $V$ on $L_2$ is the same point $L_2\cap L_6$ in all $\mathcal A_s$.
Similarly, on $L_3$ the slope is $\eta>0$, and the order by $x$-coordinate coincides with the order by $y$-coordinate. Here $V$ is the leftmost intersection point, and its neighbor is the triple point $Q=L_1\cap L_3\cap L_5$ with $y(Q)<-1$. The new intersection points have $|y|<1/15$, so they lie above $Q$, that is, to the right of it.
A triangular face with the vertex $V$ would be bounded by $L_2$, $L_3$, and a line $D$ whose intersection points with both of them are neighbors of $V$. But then $D=L_6$ and at the same time $D\in\{L_1,L_5\}$, a contradiction.

\emph{Conclusion.} No triangular face of $\mathcal A_s$ meets $R$, so all of them remain faces of $\mathcal A_s\cup\{R\}$.
Conversely, let $F$ be a triangular face of $\mathcal A_s\cup\{R\}$ without a side on $R$.
It is contained in a unique face $G$ of $\mathcal A_s$. If $R$ cut $G$, then $F$ would be one of the pieces and would have a side on $R$. Hence $R$ does not cut $G$, and $F=G$.
Thus the triangular faces of $\mathcal A_s\cup\{R\}$ are exactly the triangular faces of $\mathcal A_s$ and the $R$-triangles, which gives \eqref{eq:a3-with-R}.
\end{proof}

\subsection{\texorpdfstring{Proof of Theorem~\ref{thm:series}}{Proof of the theorem}}

We check that \cref{lem:doubling,lem:r-triangle} apply to each $\mathcal A_s$, $0\le s<t$.
For $\mathcal A_0$, where $n=6$, we check (D1) on the lines \eqref{eq:base-lines}. The $x$-intercepts of $L_1,\ldots,L_4$ are $-\sqrt3,-1/\sqrt3,1/\sqrt3,\sqrt3$, that is, $\tan(k\pi/6)$ with $1\le|k|\le2$, and for the special pair they are $-\varepsilon<0<\varepsilon$, where $\varepsilon\le1/24<1/6$ by \eqref{eq:A-num}.
Condition (D2) holds: $L_{n-1}=L_5$ (with $x$-intercept $-\varepsilon<0$) and $L_n=L_6$ have slopes $u<0<v$ by \cref{lem:uv}.
Finally, we check (D3). By \cref{lem:seed-orders}, the lines cross $Y_0$ in the order $L_1,L_2,L_5,L_6,L_3,L_4$, and by \cref{app:seed-faces} the triangular faces of $\mathcal A_0$ with a side on $Y_0$ are
\[
Y_0L_1L_2,\quad Y_0L_2L_5,\quad Y_0L_5L_6,\quad Y_0L_3L_6,\quad Y_0L_3L_4,
\]
exactly one on each bounded segment.
They alternate between the two sides of $Y_0$: by the argument at the end of the proof of \cref{lem:doubling}, two consecutive triangles on the same side would require three consecutive lines among $L_1,L_2,L_5,L_6,L_3,L_4$ to pass through one point, while the only multiple points $P=L_1\cap L_4\cap L_6$ and $Q=L_1\cap L_3\cap L_5$ are not of this form.
The transition from $\mathcal A_s$ to $\mathcal A_{s+1}$ is provided by \cref{lem:doubling}: under condition \eqref{eq:special-slopes}, the arrangement $\mathcal B$ satisfies (D1)--(D3) with $2n$ in place of $n$. At step $s$, condition \eqref{eq:special-slopes} means $\varepsilon<1/2n_s$, which holds since $\varepsilon=1/4n_t\le1/4n_s<1/2n_s$.
Therefore the original pair $L_5,L_6$ remains the special pair in all $\mathcal A_s$, so the signs $u<0<v$ and $\sigma=-1$ carry over as well.

The other hypotheses of \cref{lem:r-triangle} also hold.
First, $\delta=2n_t^6\ge2n_s^6$, since $n_s\le n_t$.
Second, the slope of the line $L_{n_s-2}$ is at most $1/2$ in absolute value and has the required sign.
For $s=0$, this line is $L_4$ with positive slope $2/5$, as required for $n=6$.
For $s\ge1$, it is the line $M_j$ of step $s-1$ with $j=n_{s-1}$ and slope $|\mu_j|<2m_{\min}(\mathcal A_{s-1})/n_{s-1}^{10}<1/2$, since $m_{\min}(\mathcal A_{s-1})\le m_{\min}(\mathcal A_0)=2/5$ by \eqref{eq:m-min} and \cref{lem:uv}.
Its parameter $\beta_j=\pi/2-\pi/2n_{s-1}$ lies in $(0,\pi/2)$, so $\sin(2\beta_j)>0$ and $b_j>0$, and by \eqref{eq:new-lines}, $\mu_j<0$.
Third, condition (W1) for $s=0$ is \cref{cor:rightmost-L4}, and the transition from $s$ to $s+1$ is provided by \cref{lem:first-step,lem:right-wedges}.
Fourth, condition (W2): for $s=0$ the strip \eqref{eq:strip} is $|y|<m_{\min}(\mathcal A_0)/n_0=1/15$ by \cref{lem:uv}. By \cref{lem:r-triangles-below}, the four $R$-triangles other than $\Delta_6=RY_0L_4$ lie in the half-plane $y\le-1$ and therefore do not meet the strip, whereas $\Delta_6$ contains the point $(\delta,0)$ and meets it. The transition from $s$ to $s+1$ is provided by \cref{lem:r-triangle}.

The arrangement $\mathcal A_0$ has exactly two triple points $P$, $Q$ and no other multiple points (\cref{lem:seed}). By \cref{lem:doubling}, each doubling step preserves the multiple points together with their multiplicities and creates no new ones, and adding $R$ does not change multiplicities (\cref{lem:r-avoids}). Hence $\mathcal C_t$ has exactly two triple points.

For $0\le s\le t$, let $A_s=a_3(\mathcal A_s)$ and $r_s=r(\mathcal A_s\cup\{R\})$. By \cref{lem:seed},
\[
A_0=10,
\qquad
r_0=5.
\]
For $0\le s<t$, \cref{lem:doubling,lem:r-triangle} give
\[
A_{s+1}=A_s+n_s^2,
\qquad
r_{s+1}=r_s+\frac{n_s}{2}.
\]
Summing the geometric progressions
\[
\sum_{s=0}^{t-1}n_s^2=36\sum_{s=0}^{t-1}4^s=12\left(4^t-1\right),
\qquad
\sum_{s=0}^{t-1}\frac{n_s}2=3\sum_{s=0}^{t-1}2^s=3\left(2^t-1\right),
\]
and substituting $n_t^2=36\cdot4^t$, $n_t=6\cdot2^t$, we obtain
\[
A_t=10+12\left(4^t-1\right)=\frac{n_t^2}{3}-2,
\qquad
r_t=5+3\left(2^t-1\right)=\frac{n_t}{2}+2,
\]
and therefore, by \cref{lem:r-avoids},
\[
a_3(\mathcal C_t)=A_t+r_t
=\frac{n_t^2}{3}+\frac{n_t}{2}.
\]
Since $N_t=n_t+2$,
\[
a_3(\mathcal C_t)=
\frac{N_t(N_t-\frac52)+1}{3},
\]
which is \eqref{eq:a3-series}. \qed
